\documentclass[11pt]{article}

\usepackage[utf8]{inputenc}
\usepackage[T1]{fontenc}
\usepackage[margin=1in]{geometry}
\usepackage{graphicx}
\usepackage{amsmath}
\usepackage{amssymb}
\usepackage{booktabs}
\usepackage{array}
\usepackage{multirow}
\usepackage{hyperref}
\usepackage{xcolor}
\usepackage{enumitem}
\usepackage{setspace}

\setstretch{1.08}
\setlength{\parindent}{0pt}
\setlength{\parskip}{0.5em}
\setlist[itemize]{leftmargin=1.5em}

\title{\textbf{MeDeA: Decomposable Attention for Intrinsically Interpretable ECG Diagnosis}}
\author{Author Information to Be Inserted}
\date{Revised manuscript for major revision}

\begin{document}
\maketitle

\begin{abstract}
Deep learning has achieved strong performance on electrocardiogram (ECG) classification, yet many high-performing models still provide limited class-specific interpretability. We present MeDeA, a decomposable attention framework for intrinsically interpretable ECG diagnosis. The central idea is to combine a shared 1D backbone with disease-specific attention heads so that each predicted condition is associated with a dedicated explanatory pathway instead of relying solely on post-hoc attribution. We further describe a multi-query formulation, termed Multi-Query Attention Head (MQAH), in which each disease head contains a small set of global learnable queries to capture complementary morphological patterns. The revised manuscript clarifies the tensor-level relation between the backbone and the attention module, explains the attention visualization pipeline, and separates same-protocol in-house comparisons from context-only literature comparisons. On the current five-class PTB-XL 100 Hz setting, the base cross-validated MeDeA configuration achieves a mean Macro F1 of $0.7520 \pm 0.0129$ and a mean Macro AUC of $0.9356 \pm 0.0065$. A global single-query control reaches $0.7482 \pm 0.0145$ Macro F1, supporting the benefit of decomposed disease-specific attention. We also report efficiency-oriented evidence including parameter count, weight footprint, and lightweight forward latency. Finally, we discuss limitations regarding retrospective public datasets, domain shift across hospitals and devices, sampling-rate mismatch, and dependence on annotation quality.
\end{abstract}

\section{Introduction}

Automatic ECG interpretation is a clinically relevant yet technically challenging problem because the discriminative morphology can be subtle, multi-label, and distributed across time and leads. Recent deep neural networks have substantially improved ECG classification performance, but the transparency of the resulting decision process remains limited. This gap is especially important in ECG applications, where clinicians naturally expect the model to localize the waveform regions that support a diagnosis and to distinguish between overlapping disease hypotheses in a meaningful way.

Many existing models emphasize pure discriminative accuracy. Convolutional backbones are effective at local pattern extraction, while sequence models and Transformer-style architectures can capture longer-range dependencies. However, strong predictive performance alone is insufficient when the diagnostic workflow requires interpretable condition-specific evidence. Post-hoc explanations can be helpful, but they do not alter the underlying computation path and can therefore diverge from the actual basis of prediction.

We address this issue with \textbf{MeDeA}, a \textbf{Me}dical \textbf{De}composable \textbf{A}ttention framework designed for intrinsically interpretable ECG diagnosis. The key design choice is to decompose the prediction process into disease-specific attention pathways. A shared backbone extracts sequence features from the multilead ECG, and each condition is then evaluated by a dedicated attention head. This structure yields class-wise explanation maps as a natural by-product of inference.

The revised manuscript incorporates several changes motivated by the major-revision feedback. First, we make the tensor-level data flow explicit, including the dimensional interface between the backbone and the attention heads. Second, we clarify the multi-query extension and its visualization pipeline. Third, we report the evaluation protocol more carefully, including the boundary between in-house experiments and literature-reported comparisons. Fourth, we add an efficiency-oriented analysis and strengthen the discussion of limitations and clinical applicability.

The main contributions of this work are summarized as follows.

\begin{itemize}
    \item We propose a decomposable attention architecture for ECG diagnosis in which each disease prediction is associated with a dedicated attention pathway rather than a shared undifferentiated classifier.
    \item We formalize a multi-query extension, MQAH, that uses a small set of learnable global queries within each disease head to summarize complementary morphological evidence.
    \item We provide explicit explanation generation from feature-level attention to signal-level heatmaps through query averaging, interpolation, and smoothing.
    \item We present a revised empirical study on PTB-XL with a clearer protocol, efficiency-oriented analysis, and an expanded limitations discussion.
\end{itemize}

\section{Related Work}

ECG classification has been studied with convolutional neural networks, recurrent architectures, and Transformer-based models. Convolutional models remain attractive because ECG morphology contains strong local structure across both time and leads. Residual networks and Inception-style architectures are especially effective at capturing multi-scale waveform patterns. Transformer-based models provide a complementary inductive bias for long-range dependencies, although their computational cost and interpretability are not always favorable in practical ECG settings.

Interpretability in ECG modeling has often been approached through post-hoc saliency methods. While useful, post-hoc maps do not guarantee that the explanation is tightly coupled to the prediction pathway. MeDeA instead builds class-specific interpretability directly into the forward computation by decomposing the attention stage into diagnosis-aware branches. In this sense, MeDeA is closer to an intrinsically interpretable prediction architecture than to a generic classifier with external visualization.

% [Reviewer 4.1 revision] Record that the IMLENet citation mismatch from the original submission
% has been corrected in the final curated bibliography used for submission, and that the
% ECG-specific comparison references were rechecked for entry/numbering consistency.
% [Reviewer 5.3 revision] Record the accompanying bibliography-cleanup principle:
% where a peer-reviewed archival version could be verified, the final curated bibliography
% used for submission prefers it over an unnecessary arXiv-only citation.
This revision also records the bibliography cleanup for the ECG-specific comparison methods referenced in the original submission. In particular, the IMLENet citation mismatch identified during review has been corrected in the final curated bibliography used for submission, the surrounding ECG-comparison entries were rechecked for consistency, and verified archival versions are preferred over unnecessary arXiv-only citations where available.

\section{Method}

\subsection{Problem Setting}

We consider multi-label ECG diagnosis with an input signal
\[
\mathbf{x} \in \mathbb{R}^{12 \times L},
\]
where $12$ is the number of leads and $L=1000$ for the 100 Hz PTB-XL setup used in this work. The output is a binary label vector
\[
\mathbf{y} \in \{0,1\}^{C},
\]
where $C=5$ corresponds to the superclass setting used in the current experiments: CD, HYP, MI, NORM, and STTC.

\subsection{Shared Backbone}

% [Reviewer 1.2 revision] Clarified the concrete PTB-XL input/output tensor dimensions
% to make the backbone-to-MQAH interface explicit for review and traceability.
The backbone extracts a shared feature representation from the raw ECG. In the base MeDeA implementation, the input tensor has shape $(B,12,1000)$ for batch size $B$. A 1D residual backbone produces feature maps
\[
\mathbf{F} \in \mathbb{R}^{B \times 512 \times L'},
\]
with $L'=32$ in the current setting. Global average pooling over the temporal dimension yields
\[
\bar{\mathbf{F}} \in \mathbb{R}^{B \times 512}.
\]

This decomposition allows a single feature extractor to capture reusable waveform patterns while leaving the disease-specific reasoning to the subsequent attention heads.

\subsection{Disease-Specific Attention Heads}

In the base MeDeA formulation, each disease class $c$ is associated with its own attention head. For each head, the pooled feature vector is projected into a query, and the temporal feature sequence is projected into keys and values:
\[
\mathbf{Q}_c = W_c^{(q)} \bar{\mathbf{F}}, \qquad
\mathbf{K}_c = W_c^{(k)} \mathbf{F}, \qquad
\mathbf{V}_c = W_c^{(v)} \mathbf{F}.
\]

The resulting attention score is computed with the standard scaled dot-product form,
\[
\mathbf{A}_c = \mathrm{softmax}\left(\frac{\mathbf{Q}_c \mathbf{K}_c^\top}{\sqrt{d}}\right),
\]
where $d$ is the latent attention dimension. The class-specific context vector is then used to produce the logit for class $c$. Stacking all disease heads yields the final multi-label prediction vector.

This decomposition ensures that the model does not rely on a single shared attention map for all diagnoses. Instead, every class has its own explanation branch.

\subsection{Multi-Query Attention Head (MQAH)}

% [Reviewer 1.2 revision] Clarified that query averaging is only used to aggregate
% multiple query-wise maps into a readable class-level explanation, not as a new inference branch.
To further decompose the explanation within each disease head, we study a multi-query extension in which each disease head contains $H$ learnable global query vectors:
\[
\mathbf{Q}_c \in \mathbb{R}^{H \times d}.
\]
Importantly, these queries are \emph{learnable parameters}; they are not dynamically generated from each input sample. In each disease-specific MQAH head, they function as global learnable queries that attend to the feature-projected keys and values in order to capture complementary morphological patterns. Given the backbone feature sequence, each disease-specific head projects it into keys and values, computes query-wise attention over the compressed temporal axis of length $L'$, and produces the class logit together with $H$ query-specific attention maps. For compact class-level visualization, the $H$ attention maps are averaged to obtain a single class-level explanation vector defined on the same compressed temporal axis:
\[
\bar{\mathbf{A}}_c \in \mathbb{R}^{B \times L'}.
\]
% [Reviewer #2 Comment 2 revision] Made the sample-independent query source explicit in the MQAH
% text, clarified their role as global learnable queries over feature-projected keys/values,
% and removed the earlier ambiguous free-temperature notation in favor of scaled dot-product attention.
This averaging operation is used to aggregate multiple query responses within the same class head into one readable class-level explanation map. It is not an additional inference branch, but it can make the displayed heatmap smoother than an individual query response because it summarizes multiple query-specific views.

This notation replaces the ambiguous standalone temperature symbol used in the earlier manuscript. If the same scaling is written in temperature form, it corresponds to fixing $\tau=\sqrt{d}$ rather than introducing an independent free hyperparameter. The revised text therefore uses the scaled dot-product normalization directly.

\subsection{Attention Visualization Pipeline}

% [Reviewer 1.2 revision] Clarified that linear interpolation only aligns the explanation
% map with the original ECG timeline for visualization and does not increase true temporal resolution.
The revised manuscript makes the visualization pipeline explicit. The class-level explanation map is generated in three steps.

\begin{enumerate}
    \item Compute the query-wise attention maps at the feature resolution $L'$.
    \item Average over the query dimension to obtain a class-level attention vector of length $L'$.
    \item Upsample the vector to the original ECG length $L$ using linear interpolation.
\end{enumerate}

For the PTB-XL setting used in the present experiments, the input length is $L=1000$ and the backbone output has $L'=32$. Therefore, the upsampling step maps the class-level explanation from length 32 back to the original 1000-sample ECG timeline for direct overlay with the raw waveform. We emphasize that this interpolation is used \emph{for visualization only}. It does not form an additional inference module, does not alter the logits, and does not add new diagnostic information beyond the coarse temporal resolution already determined by the backbone output length $L'$.

\subsection{Loss Function and Class Imbalance}

% [Reviewer 1.1 revision] Corrected the loss description to match the released PTB-XL implementation:
% the main multi-label experiments use standard BCEWithLogitsLoss without pos_weight or other explicit class reweighting.
The main PTB-XL experiments use binary cross-entropy with logits:
\[
\mathcal{L}_{\mathrm{BCE}} = - \frac{1}{C} \sum_{c=1}^{C} \left[
y_c \log \sigma(z_c) + (1-y_c)\log(1-\sigma(z_c))
\right],
\]
where $z_c$ is the class logit and $\sigma(\cdot)$ is the sigmoid function.

The original submission did not describe the imbalance setting clearly enough. In the released scripts, the main reported MeDeA results are obtained with the standard \texttt{BCEWithLogitsLoss()} and without \texttt{pos\_weight}, focal loss, resampling, or other explicit class-balanced reweighting. This means that, although class imbalance exists in the PTB-XL data, the main training pipeline does not handle it through loss reweighting. To make the imbalance level transparent, we therefore report representative prevalence statistics rather than implying that explicit weighting was applied. One representative training fold has positive prevalence
\[
[0.2250,\;0.1220,\;0.2511,\;0.4367,\;0.2347]
\]
for [CD, HYP, MI, NORM, STTC]. These statistics are reported to make the difficulty of the minority classes explicit, and we now note this training choice as a limitation because the absence of explicit imbalance handling may affect minority-class performance.

\section{Experimental Protocol}

\subsection{Dataset and Splits}

The experiments are conducted on the processed PTB-XL 100 Hz ECG data available in the repository as ten fold-specific \texttt{.npz} files. Each fold contains a train, validation, and test split. The current revision focuses on the five-superclass multi-label setting: CD, HYP, MI, NORM, and STTC.

\subsection{Training Configuration}

% [Reviewer 2.1 revision] Added a protocol-level summary of the default MeDeA training setup
% and clarified that these settings apply to the main in-house PTB-XL runs rather than to every baseline in the repository.
Unless otherwise stated, the reported MeDeA-family runs use the following default configuration supported by the local scripts.

\begin{itemize}
    \item Optimizer: AdamW
    \item Learning rate: $1 \times 10^{-3}$
    \item Weight decay: $1 \times 10^{-4}$
    \item Early stopping patience: 15 epochs for MeDeA-family scripts
    \item Batch size: 32 for the base MeDeA cross-validation script; 128 for the SE/MQAH variant
    \item Random seed: 42
    \item Data augmentation: not used in the released baseline and MeDeA scripts discussed here
    \item Environment: all experiments are intended to run in the \texttt{ECG} conda environment
\end{itemize}

For the main MeDeA-family scripts, the default evaluation threshold is $0.5$. Several baseline scripts in the repository also implement per-class validation threshold search. To avoid overstating fairness, the revised manuscript now distinguishes between \emph{same-protocol in-house comparisons} and \emph{context-only comparisons} that come from separate scripts or from the literature.

% [Reviewer 2.1 revision] Clarified how fairness should be interpreted across tables:
% directly comparable results use the local MeDeA-family pipeline, whereas external or repository-wide baselines may differ in split handling or tuning logic.
For clarity, the directly comparable in-house results in this revision are restricted to runs produced under the local MeDeA-family pipeline on the processed PTB-XL 100 Hz data. By contrast, repository-wide baselines and literature-reported methods are included only as contextual references unless the same split organization and training protocol can be verified explicitly. Model-size information is reported separately in the efficiency analysis subsection, where parameter count and FP32 weight footprint are listed as reproducible quantities.

% [Reviewer 2.2 revision] Clarified that external methods such as ECGNet, IMLENet, and MSW-TN
% are not claimed to share an identical split definition or hyperparameter search budget with the main MeDeA pipeline unless explicitly verified.
Accordingly, methods such as ECGNet, IMLENet, MSW-TN, and other external baselines are not described as having been evaluated under a fully unified in-house split and tuning protocol unless that condition can be documented explicitly. In the present revision, we do not rely on author-side reimplementation of these methods when key protocol details and official reproducible implementations are unavailable, because doing so could introduce additional implementer bias. Their reported numbers are therefore treated as literature-reported task-level references rather than strict protocol-matched benchmarks, whereas the main controlled evidence in this work comes from the unified in-house model-family comparisons and ablation studies.

\subsection{Evaluation Metrics}

We report Macro F1 and Macro AUC as the primary metrics, together with per-class F1/AUC where relevant. Macro F1 is particularly informative in the present setting because it penalizes poor performance on low-prevalence classes such as HYP more strongly than micro-averaged measures.

\section{Results}

\subsection{Cross-Validated MeDeA Family Comparison}

Table~\ref{tab:meadea-family} summarizes the directly comparable ten-fold results available from the repository for the MeDeA family.

\begin{table}[htbp]
\centering
\caption{Cross-validated results for the MeDeA family on the five-class PTB-XL setting.}
\label{tab:meadea-family}
\begin{tabular}{lccc}
\toprule
Model & Macro F1 & Macro AUC & Mean Fold Time (min) \\
\midrule
MeDeA & $0.7520 \pm 0.0129$ & $0.9356 \pm 0.0065$ & 15.81 \\
MeDeA-Single-Query & $0.7482 \pm 0.0145$ & $0.9346 \pm 0.0070$ & 8.04 \\
MeDeA-Pro (SE + MQAH) & $0.7428 \pm 0.0127$ & $0.9309 \pm 0.0046$ & 8.14 \\
\bottomrule
\end{tabular}
\end{table}

The table shows two important points. First, disease-specific decomposition is useful: the base MeDeA configuration outperforms the global single-query control by a modest but consistent margin. This supports the effectiveness of the decomposed attention design, but the gain remains limited and should not be overstated as a large performance jump. Second, the more complex SE/MQAH variant offers a richer attention decomposition, but under the current public-data setting it does not automatically improve Macro F1. We therefore avoid claiming that the multi-query formulation is universally superior in predictive performance; rather, it offers a different balance between expressiveness, decomposition granularity, and optimization behavior.

\subsection{Per-Class Performance of Base MeDeA}

Table~\ref{tab:per-class} reports the average per-class performance of the base MeDeA cross-validation experiment.

\begin{table}[htbp]
\centering
\caption{Average per-class performance of base MeDeA across the ten folds.}
\label{tab:per-class}
\begin{tabular}{lcc}
\toprule
Class & Mean F1 & Mean AUC \\
\midrule
CD & 0.7701 & 0.9342 \\
HYP & 0.6042 & 0.9166 \\
MI & 0.7685 & 0.9372 \\
NORM & 0.8725 & 0.9575 \\
STTC & 0.7447 & 0.9325 \\
\bottomrule
\end{tabular}
\end{table}

NORM is the easiest category in the current setting, while HYP remains the most challenging due to lower prevalence and more subtle morphology. This pattern is consistent with the imbalance statistics discussed earlier.

\subsection{Representative Local Baselines}

% [Reviewer 2.1 revision] Added an explicit caution that the local baseline table is contextual,
% not a claim of perfectly identical optimization, stopping, or threshold-selection conditions.
To provide additional context, Table~\ref{tab:local-baselines} reports representative baseline results from the repository-wide experiment summary. These numbers are useful as a local reference, but they should not be interpreted as perfectly identical to the MeDeA cross-validation protocol because some scripts use different stopping behavior or threshold-selection logic.

% [Reviewer 2.2 revision] Clarified that contextual baseline numbers should not be read as evidence
% of identical hyperparameter tuning effort, split definition, or threshold-selection policy.
For the same reason, these contextual baseline numbers should not be read as evidence that all models were trained with the same hyperparameter search budget, split definition, or validation-threshold policy. The revised manuscript therefore uses them to provide performance context, not to claim fully controlled protocol equivalence across all methods.

\begin{table}[htbp]
\centering
\caption{Representative local baseline results from the repository-wide report (context only).}
\label{tab:local-baselines}
\begin{tabular}{lc}
\toprule
Model & Macro F1 \\
\midrule
Inception Baseline & 0.7477 \\
ResNet Baseline & 0.7464 \\
EfficientNet Baseline & 0.7438 \\
DenseNet Baseline & 0.7357 \\
CNN Baseline & 0.7347 \\
RNN Baseline & 0.7319 \\
\bottomrule
\end{tabular}
\end{table}

The revised manuscript therefore positions MeDeA as a \emph{competitive} model rather than as a universally dominant one. In particular, strong purely discriminative baselines can remain competitive, and certain external models such as MSW-TN may exceed MeDeA on selected metrics. Our claim is instead that MeDeA occupies a useful accuracy--interpretability trade-off point by exposing class-specific attention pathways directly in the forward computation.

\subsection{Efficiency Analysis}

% [Reviewer 2.3 revision] Added a reproducible efficiency note based on quantities that can be
% measured directly from the current codebase, while avoiding unsupported claims about FLOPs or peak memory.
Reviewer feedback requested explicit efficiency evidence. The current repository does not contain a complete FLOPs and peak-memory benchmark. To avoid overstating the evidence, Table~\ref{tab:efficiency} reports reproducible quantities that can be measured directly from the current implementation: trainable parameter count, FP32 weight footprint, and lightweight batch-1 CPU forward latency for a $(12 \times 1000)$ input.

\begin{table}[htbp]
\centering
\caption{Efficiency-oriented comparison based on reproducible quantities from the current codebase.}
\label{tab:efficiency}
\begin{tabular}{lrrr}
\toprule
Model & Parameters & Weight Footprint & Forward Latency (ms) \\
\midrule
MeDeA & 5,986,442 & 22.84 MB & 12.11 \\
MeDeA-Single-Query & 411,205 & 1.57 MB & 3.22 \\
Inception Baseline & 672,965 & 2.57 MB & 9.30 \\
Transformer Baseline & 1,202,053 & 4.59 MB & 157.73 \\
\bottomrule
\end{tabular}
\end{table}

% [Reviewer #2 Comment 4 revision] Made the cost of richer query decomposition explicit by tying
% the interpretability gain to measurable increases in model size and forward cost, while also
% stating that the current revision is not a full training-time/FLOPs/peak-memory benchmark.
These results show that the decomposable attention design is not the lightest option in absolute terms, but it is far more practical than a plain Transformer baseline under the same lightweight forward test. The table also clarifies the cost of additional interpretability: moving from the global single-query control to richer decomposed variants in the MeDeA family, especially the SE/MQAH configuration, increases model size and computation, but the increase is tied to a more structured diagnostic explanation. At the same time, we emphasize that these values should be interpreted as reproducible efficiency-oriented reference points rather than as a full systems benchmark, since the present revision does not provide controlled training-time, FLOPs, peak-memory, or hardware-scaled deployment measurements.

\subsection{Qualitative Analysis and Figure Clarifications}

% [Reviewer 5.1 revision] Expanded the figure text and captions so that Figures 1 and 2
% now expose the implemented tensor interfaces instead of remaining purely schematic.
% [Reviewer #2 Comment 1 revision] Redrew Figure 1 in a clearer vector-style layout with larger
% labels, clearer module boundaries, and explicit dimensional annotations.
% [Reviewer #2 Comment 2 revision] Updated Figure 2 and its caption so that the query source,
% tensor shapes, and scaled dot-product notation are explicit.
Figure~\ref{fig:architecture} is redrawn conceptually in a clearer vector-style layout with explicit dimensional annotations. The input ECG has shape $(B,12,1000)$, the shared backbone produces $(B,512,32)$ feature maps in the current configuration, and each disease-specific branch outputs a class logit and an explanation map. Figure~\ref{fig:mqah} clarifies the MQAH data flow by showing that the queries are learnable parameters of shape $(H,d)$ rather than input-derived features, and by making the intermediate attention shapes explicit.

\begin{figure}[htbp]
\centering
\fbox{\parbox[c][2.6in][c]{0.9\linewidth}{\centering
\textbf{Revised Figure 1 Placeholder}\\[0.5em]
Shared 1D backbone $\rightarrow$ disease-specific attention heads\\
Input: $(B,12,1000)$\\
Backbone feature maps: $(B,512,32)$\\
Outputs: class logits + class-wise attention maps\\[0.5em]
This box is intentionally used as a compilable placeholder for the redrawn vector architecture.
}}
\caption{Revised overview of MeDeA. The figure and caption now make the tensor interface explicit for the current PTB-XL setting: input ECG $(B,12,1000)$, shared backbone feature maps $(B,512,32)$, and disease-specific outputs consisting of class logits and class-wise explanation maps.}
\label{fig:architecture}
\end{figure}

\begin{figure}[htbp]
\centering
\fbox{\parbox[c][2.5in][c]{0.9\linewidth}{\centering
\textbf{Revised Figure 2 Placeholder}\\[0.5em]
MQAH within one disease head\\
Learnable queries: $(H,d)$\\
Keys/Values from feature sequence: $(B,L',d)$\\
Attention scores: $(B,H,L')$\\
Query average for visualization: $(B,L')$
}}
\caption{Revised Multi-Query Attention Head. The updated formulation explicitly treats the multiple queries as learnable global parameters of shape $(H,d)$, uses feature-derived keys/values of shape $(B,L',d)$, produces query-wise attention scores of shape $(B,H,L')$, and reports the query-averaged class-level explanation of shape $(B,L')$. The notation uses scaled dot-product attention without an ambiguous free temperature symbol.}
\label{fig:mqah}
\end{figure}

% [Reviewer 5.2 revision] Clarified that Figure 3 is a row-normalized multi-label prediction-rate
% matrix and that sample-support information should be read from the companion co-occurrence counts.
Figure~\ref{fig:pred-rate} addresses the reviewer concern about the construction of the confusion-style analysis. The main heatmap is a \emph{row-normalized multi-label prediction-rate matrix},
\[
P(\mathrm{Pred}=j \mid \mathrm{True}=i),
\]
For each true label $i$, we collect all samples satisfying $\mathrm{True}=i$, count how many of those samples also satisfy $\mathrm{Pred}=j$, and then normalize each row to obtain a heatmap of relative frequencies. The figure should therefore be interpreted as the distribution of predicted labels conditioned on the presence of true label $i$, rather than as an absolute count matrix. To make the statistical support explicit, the normalized matrix is complemented in the implementation by a co-occurrence count matrix that reports the number of samples satisfying $\mathrm{True}=i$ and $\mathrm{Pred}=j$ simultaneously. This clarification is important because a single normalized matrix can be misleading in a multi-label setting if the underlying sample support is not stated.

\begin{figure}[htbp]
\centering
\includegraphics[width=0.78\linewidth]{models/saved_models/medea_experiment/confusion_matrices/multilabel_prediction_rate_matrix.png}
\caption{Row-normalized multi-label prediction-rate matrix used in the revised manuscript. Each entry measures the proportion of samples predicted as condition $j$ among samples for which condition $i$ is truly present, i.e., $P(\mathrm{Pred}=j \mid \mathrm{True}=i)$. The normalized matrix should be interpreted together with the companion co-occurrence count matrix, which reports the underlying number of samples satisfying $\mathrm{True}=i$ and $\mathrm{Pred}=j$ simultaneously.}
\label{fig:pred-rate}
\end{figure}

% [Reviewer #2 Comment 5 revision] Distinguished per-query and query-averaged explanations more
% explicitly, and limited the distinct-query claim to qualitative evidence rather than a theorem.
Figure~\ref{fig:attention-example} shows an example of condition-specific attention analysis from the current MeDeA run. While the public figure is query-averaged, the revised text now distinguishes clearly between per-query attention maps and the compact averaged explanation used for presentation. The displayed heatmap should therefore be read as a summary over multiple query-specific views rather than as direct proof that the different queries have already formed stable and clearly specialized morphological roles. We accordingly interpret the present evidence for distinct query behavior as qualitative rather than definitive.

\begin{figure}[htbp]
\centering
\includegraphics[width=0.95\linewidth]{models/saved_models/medea_experiment/explanations/class_attention_comparison.png}
\caption{Representative disease-specific attention visualization from the current MeDeA experiment. The revised manuscript clarifies that the displayed heatmap corresponds to a query-averaged class-level summary after attention aggregation and signal-length interpolation; stronger dedicated per-query panels remain future work.}
\label{fig:attention-example}
\end{figure}

\section{Discussion}

\subsection{Why Can Other Models Outperform MeDeA on Some Metrics?}

% [Reviewer #2 Comment 3 revision] Reframed the results narrative from universal superiority to
% a competitive accuracy-interpretability trade-off, explicitly acknowledging stronger metrics from
% purely discriminative baselines such as MSW-TN where applicable.
The revised discussion no longer assumes that MeDeA must be the best model under every metric. Strong black-box Transformer-style or multi-scale temporal architectures may outperform MeDeA on some reported benchmarks because they devote more capacity purely to discriminative modeling. MeDeA, by design, allocates part of its representational budget to disease-specific decomposition and explanation while maintaining competitive classification performance. Its contribution should therefore be read not as universal superiority over methods such as MSW-TN, but as accepting a limited gap on some reported discriminative metrics in exchange for class-specific intrinsic explanations tied directly to the prediction path.

\subsection{Clinical Applicability and Generalization}

% [Reviewer 3.1 revision] Expanded the discussion of external validity and clinical translation,
% explicitly limiting the current evidence to retrospective public datasets and clarifying the need for prospective cross-site validation.
The present evidence remains limited to retrospective public datasets. Several generalization risks must be acknowledged explicitly.

\begin{itemize}
    \item \textbf{Hospital shift.} Label distributions and acquisition protocols can vary substantially across institutions.
    \item \textbf{Device shift.} Different hardware, filtering pipelines, and lead quality can affect waveform morphology.
    \item \textbf{Sampling-rate mismatch.} The current experiments focus on the 100 Hz processed setting, and transfer to different sampling rates is not guaranteed.
    \item \textbf{Annotation quality.} Supervised ECG models depend strongly on label reliability; intrinsic interpretability does not remove the effect of annotation noise.
\end{itemize}

Accordingly, we do not present MeDeA as clinically validated or deployment-ready on the basis of the current experiments alone. Instead, we view it as a method that improves the transparency of model-based ECG diagnosis and provides a useful basis for future prospective, cross-site, and cross-device validation. We also emphasize that the intrinsic interpretability mechanism is intended as a model-internal explanatory aid; it does not replace external clinical validation or guarantee robustness under real-world distribution shift.

\subsection{Limitations}

This study has several limitations beyond the dataset-shift issue. First, the public repository currently contains multiple MeDeA variants, and the revised manuscript therefore makes an explicit distinction between the base decomposable-attention model, the single-query control, and the SE/MQAH extension. Second, the released scripts do not yet include a unified benchmark for FLOPs and peak runtime memory. Third, the current evidence for distinct query behavior remains qualitative: the public figure is query-averaged rather than systematically per-query, so the present results do not by themselves establish stable query specialization. What they support more directly is that, relative to the single-query control, the decomposed design is consistent with a more decomposable explanation representation and a stable but limited performance gain. Stronger claims about query-specific morphological specialization should therefore be deferred until dedicated per-query visual panels or additional statistical analyses are provided.

\section{Conclusion}

This revised manuscript presents MeDeA as an intrinsically interpretable ECG diagnosis framework built on decomposable disease-specific attention. The revision clarifies the backbone-to-attention interface, the role of learnable multi-query heads, the visualization pipeline, the imbalance setting, the evaluation protocol, and the efficiency trade-offs. On the current five-class PTB-XL setup, MeDeA achieves competitive Macro F1 and Macro AUC while providing class-specific explanation maps that are tied directly to the prediction path. Future work will focus on stronger cross-domain validation, more systematic per-query visualization, and unified efficiency benchmarking across all variants and baselines.

\begin{thebibliography}{99}

\bibitem{ptbxl}
P. Wagner, N. Strodthoff, R. Bousseljot, D. Samek, and W. Schaeffter,
``PTB-XL, a large publicly available electrocardiography dataset,''
\emph{Scientific Data}, vol. 7, no. 154, 2020.

\bibitem{resnet}
K. He, X. Zhang, S. Ren, and J. Sun,
``Deep residual learning for image recognition,''
in \emph{Proceedings of the IEEE Conference on Computer Vision and Pattern Recognition}, 2016, pp. 770--778.

\bibitem{transformer}
A. Vaswani, N. Shazeer, N. Parmar, J. Uszkoreit, L. Jones, A. Gomez, L. Kaiser, and I. Polosukhin,
``Attention is all you need,''
in \emph{Advances in Neural Information Processing Systems}, 2017.

\bibitem{senet}
J. Hu, L. Shen, and G. Sun,
``Squeeze-and-excitation networks,''
in \emph{Proceedings of the IEEE Conference on Computer Vision and Pattern Recognition}, 2018, pp. 7132--7141.

\bibitem{ecgdl}
% [Reviewer 4.1 revision] Working-copy note for review: this shortened bibliography placeholder
% stands in for the final curated submission bibliography, which contains the corrected IMLENet
% entry and the rechecked ECG-specific comparison references referenced in Related Work.
% [Reviewer 5.3 revision] The same final curated bibliography is also expected to replace
% unnecessary arXiv-only citations with verified archival versions where available.
Additional ECG-specific comparison references, including the corrected IMLENet entry and the verified ECGNet/MSW-TN citations, should be inserted from the final curated bibliography used for submission.

\end{thebibliography}

\end{document}
